\documentclass[leqno]{article}
\usepackage{verbatim}
\usepackage{array}
\usepackage{listings}
\usepackage{fancyvrb}
\usepackage{enumitem}

\usepackage{lmodern}

\usepackage[utf8]{inputenc}
\usepackage[T1]{fontenc}
\usepackage{textcomp}
\usepackage{multicol} \usepackage{mathtools}
\usepackage{amsmath}
\usepackage{wrapfig}
\usepackage{amssymb}
\usepackage{extpfeil}
\usepackage{amsmath,amsfonts,amssymb,amsthm,epsfig,epstopdf,titling,url,array}
\usepackage{hyperref}
\usepackage{eso-pic}
\usepackage{pgf}
\usepackage{tikz}
\usepackage{tikz-cd}
\usepackage{graphicx}

% figure support
\usepackage{import}
\usepackage{xifthen}
\pdfminorversion=7
\usepackage{pdfpages}
\usepackage{transparent}
\usepackage{xcolor}

% geometry
\usepackage{geometry}
\geometry{a4paper, margin=1in}

% paragraph length
\setlength{\parindent}{0em}
\setlength{\parskip}{1em}

\newtheorem*{theorem}{Theorem}
\newtheorem*{lemma}{Lemma}
\newtheorem*{proposition}{Proposition}
\newtheorem*{definition}{Definition}
\newtheorem*{observation}{Observation}

\DeclareMathOperator{\Hom}{Hom}
\DeclareMathOperator{\Poset}{Poset}
\DeclareMathOperator{\im}{Im}
\DeclareMathOperator{\coker}{Coker}
\DeclareMathOperator{\coim}{Coim}

\newcommand{\com}[1]{\textcolor{red}{#1}}
\newcommand{\incfig}[1]{%
\center
\def\svgwidth{0.9\columnwidth}
\import{./figures/}{#1.pdf_tex}
}
\newcommand{\incimg}[1]{%
\center
\includegraphics[width=0.9\columnwidth]{images/#1}
}
\pdfsuppresswarningpagegroup=1

\title{DefProp TopAlg}
\author{Abel Doñate Muñoz}
\date{}

\begin{document}
\maketitle
\tableofcontents
\newpage

\section{Teoria de grupos abelianos}

\section{Homology}
\begin{definition}[Exact sequence] Given the complex $C_{\bullet}$ formed by abelian groups $(C_n)$  and morphisms $f_n$ such that $f_{n}\circ f_{n+1} = 0$ 
    % https://q.uiver.app/#q=WzAsNSxbMCwwLCJcXGNkb3RzIl0sWzEsMCwiQ197bisxfSJdLFsyLDAsIkNfbiJdLFszLDAsIkNfe24tMX0iXSxbNCwwLCJcXGNkb3RzIl0sWzAsMSwiZl97bisyfSJdLFsxLDIsImZfe24rMX0iXSxbMiwzLCJmX3tufSJdLFszLDQsImZfe24tMX0iXV0=
\[\begin{tikzcd}
	\cdots & {C_{n+1}} & {C_n} & {C_{n-1}} & \cdots
	\arrow["{f_{n+2}}", from=1-1, to=1-2]
	\arrow["{f_{n+1}}", from=1-2, to=1-3]
	\arrow["{f_{n}}", from=1-3, to=1-4]
	\arrow["{f_{n-1}}", from=1-4, to=1-5]
\end{tikzcd}\]
the complex is an exact sequence if $\ker(f_{n}) = \im(f_n) \ \forall  n$.

If we have only three object such that $0 \to A \hookrightarrow B \twoheadrightarrow B \to 0$, and the image of the first morhphism is the kernel of the second, then we call it an short exact sequence.
\end{definition}

\begin{definition}[Homology] There exists a covariant functor from the category of chain complex to the category of chain complexes defined as
\[
H_n (C_{\bullet}) = \ker (d_n) / \im (d_{n+1})
\]
\end{definition}

\begin{proposition} If $C_{\bullet}$ is exact, then $H_{\bullet}(C_{\bullet})=0$ .
\end{proposition}

\begin{definition}[Morphism of complex] \com{palo}

\end{definition}

\begin{definition}[Chain homotopy] Two morphism $ f_{\bullet}, g_{\bullet}: C_{\bullet} \to  B_{\bullet}$ are homotopically equivalent if there exists a $h_{\bullet}$ such that $f-g = d_B\circ h + h \circ d_A$ 
% https://q.uiver.app/#q=WzAsMTEsWzEsMSwiQ157bi0xfSJdLFsyLDEsIkNee259Il0sWzMsMSwiQ157bisxfSJdLFsxLDIsIkJee24tMX0iXSxbMiwyLCJCXntufSJdLFszLDIsIkJee24rMX0iXSxbNCwxLCJcXGNkb3RzIl0sWzQsMiwiXFxjZG90cyJdLFswLDEsIlxcY2RvdHMiXSxbMCwyLCJcXGNkb3RzIl0sWzEsMF0sWzAsMSwiZF9DIl0sWzEsMiwiZF9DIl0sWzQsNSwiZF9CIl0sWzMsNCwiZF9CIl0sWzEsNF0sWzIsNSwiZiJdLFsyLDZdLFs1LDddLFs4LDBdLFs5LDNdLFsxLDQsImYiXSxbMSw0LCJnIiwyXSxbMCwzLCJmIl0sWzAsMywiZyIsMl0sWzIsNSwiZyIsMl0sWzEsMywiaF9uIiwyXSxbMiw0LCJoX3tuKzF9IiwyXV0=
\[\begin{tikzcd}
	& {} \\
	\cdots & {C^{n-1}} & {C^{n}} & {C^{n+1}} & \cdots \\
	\cdots & {B^{n-1}} & {B^{n}} & {B^{n+1}} & \cdots
	\arrow[from=2-1, to=2-2]
	\arrow["{d_C}", from=2-2, to=2-3]
	\arrow["f", from=2-2, to=3-2]
	\arrow["g"', from=2-2, to=3-2]
	\arrow["{d_C}", from=2-3, to=2-4]
	\arrow["{h_n}"', from=2-3, to=3-2]
	\arrow[from=2-3, to=3-3]
	\arrow["f", from=2-3, to=3-3]
	\arrow["g"', from=2-3, to=3-3]
	\arrow[from=2-4, to=2-5]
	\arrow["{h_{n+1}}"', from=2-4, to=3-3]
	\arrow["f", from=2-4, to=3-4]
	\arrow["g"', from=2-4, to=3-4]
	\arrow[from=3-1, to=3-2]
	\arrow["{d_B}", from=3-2, to=3-3]
	\arrow["{d_B}", from=3-3, to=3-4]
	\arrow[from=3-4, to=3-5]
\end{tikzcd}\]
\end{definition}

\begin{definition}[Chain homotopy] If
\[
\begin{cases}
f_{\bullet}: C_{\bullet} \to B_{\bullet}\\
g_{\bullet}: B_{\bullet} \to C_{\bullet}
\end{cases} 
\begin{cases}
f_{\bullet} \circ g_{\bullet} \simeq id_{B_{\bullet}}\\
g_{\bullet} \circ f_{\bullet} \simeq id_{C_{\bullet}}
\end{cases} \implies C_{\bullet} \simeq B_{\bullet}
\]

\end{definition}

\begin{theorem}[Connecting morphism] Given the short exact sequence $0 \to C _{\bullet} \xhookrightarrow{f}  D _{\bullet} \xtwoheadrightarrow{g_{\bullet}}  \to 0$ 

\end{theorem}


% ######################################################################

\section{Homology}

\begin{theorem}[Homotopical invariance] if $X \simeq Y$, then $H _{\bullet}(X) \cong H_{\bullet}(Y)$.
\end{theorem}

\begin{proposition}
We have the short exact sequence
\[
0 \to C_{\bullet}(A) \hookrightarrow C _{\bullet}(X) \twoheadrightarrow C_{\bullet}(X,A) \to 0
\]
\end{proposition}

\begin{definition}[Relative homology] $H _{\bullet}(X, A) := H( \frac{C_{\bullet}(X)}{C_{\bullet}(A)})$. This gives us the long exact sequence
\[
\cdots \to H_{n+1}(X, A) \to H_n(A) \to H_n(X) \to H_n(X,A) \to \cdots 
\]
\end{definition}

\begin{theorem}[Excision] $U\subseteq A\subseteq X$ such that $\overline{U} \subseteq A^\circ$. Then the inclusion $\iota : (X-U, A-U) \hookrightarrow (X,A)$ induces an isomorphism 
	\[
	\iota _{\bullet} :  H_n (X-U, A-U) \xrightarrow{\cong} H_n(X,A)  
	\]
\end{theorem}

\begin{definition}[Good pairs] Let $A\subseteq X$ closed. $(X,A)$ is a good pair if $A$ is a deformation retract of an open neighborhood of itself. 
\end{definition}

\begin{theorem} If $(X, A)$ is a good pair, then $q: (X, A) \twoheadrightarrow  (X/A, A/A)$ induces a isomorphism
	\[
	q_{\bullet}: H_n(X, A) \xrightarrow{\cong} H_n(X / A, A/A)  \cong \tilde{H}_n(X/A)
	\]
\end{theorem}

\begin{theorem}[Mayer-Vietoris] Let $X = X_1^\circ \cup X_2^\circ$. Then $\iota : (X_1, X_1 \cap X_2) \hookrightarrow (X, X_2)$ induces the isomorphism
	\[
	i _{\bullet}: H_n(X_1, X_1 \cap X_2) \xrightarrow{\cong}  H_n(X, X_2)
	\]
and we have the long exact sequence
\[
\cdots \to H_{n} (X_1 \cap X_2) \to H_n(X_1) \oplus H_n(X_2) \to H_n(X) \to H_{n-1}(X_1 \cap X_2) \to \cdots 
\]
\end{theorem}



% ######################################################################


\section{Fundamental group}
\com{falta}


% ######################################################################

\section{Covering spaces}

\begin{definition}[Covering space] $p: \tilde{X} \twoheadrightarrow  X$ is covering space if for every point $x \in X$ there is a neighborhood $U$ such that
\[
p^{-1}(U) = \sqcup V_j , \quad p|_{V_j}: V_j \xrightarrow{\cong} U  
\]
\end{definition}

\begin{proposition} If $p: \tilde{X} \twoheadrightarrow  X$, then
	\begin{enumerate}[topsep=-6pt, itemsep=0pt]
	\item $p$ is an open application.
	\item All the fibers have the same cardinality
	\end{enumerate}
\end{proposition}

\begin{definition}[Group acting on a space] The group $G$ acts on $Y$ if for every $g \in G$ we have an aplication $g: Y \to Y$ such that $y \mapsto gy$ and
	\begin{enumerate}[topsep=-6pt, itemsep=0pt]
		\item $1y=y \ \forall  y \in Y$
		\item $g(hy)=(gh)y \ \forall  g, h \in G, \ \forall  y \in Y$ 
	\end{enumerate}
\end{definition}

\begin{definition}[Proporly discontinous] The action of $G$ over $Y$ is properly discontinous if $\ \forall  y \in Y \ \exists U$ neighbourhood dishoint to $gU \ \forall  g \neq 1$.
\end{definition}

\begin{theorem} If $G$ acts properly discontinuos on $Y$, then $p : Y \twoheadrightarrow Y/G$ is a covering space.
\end{theorem}

\begin{theorem}[Path lifting] Given a covering space $p: \tilde{X} \twoheadrightarrow X$ and a path $\omega: I \to X$ with origin $x_0$. Fixing $\tilde{x}_0 \in p^{-1}(x_0)$, there exist a unique path $\tilde{\omega}$ such that
% https://q.uiver.app/#q=WzAsMyxbMCwxLCJJIl0sWzEsMCwiXFx0aWxkZXtYfSJdLFsxLDEsIlgiXSxbMCwyLCJcXG9tZWdhIl0sWzEsMiwicCIsMl0sWzAsMSwiXFx0aWxkZXtcXG9tZWdhfSJdXQ==
\[\begin{tikzcd}
	& {\tilde{X}} \\
	I & X
	\arrow["p"', from=1-2, to=2-2]
	\arrow["{\tilde{\omega}}", from=2-1, to=1-2]
	\arrow["\omega", from=2-1, to=2-2]
\end{tikzcd}\]
\end{theorem}

\begin{proposition} If $p: \tilde{X} \twoheadrightarrow X$ covering space, $\tilde{x}_0 \in \tilde{X}$ and $x_0\in p(\tilde{x}_0)$, then 
	\[
    p_* : \pi (\tilde{X}, \tilde{x}_0) \hookrightarrow \pi (X, x_0)
	\]
	is a monomorphism
\end{proposition}

\begin{definition}[Universal covering] $p: \tilde{X} \twoheadrightarrow X$ is a universal covering if $\tilde{X}$ is simply connected. 
\end{definition}

\begin{definition}[Regular] $p: \tilde{X} \twoheadrightarrow X$ is regular if the subgrup $p_* \pi (\tilde{X}, \tilde{x}_0) \subseteq \pi (X, x_0)$ is normal.
\end{definition}

\begin{theorem} If $p: \tilde{X} \twoheadrightarrow X$ universal covering, then $\varphi : \pi (X, x_0) \xrightarrow{\cong} p^{-1}(x_0) $ is an isomorphism.
\end{theorem}

\begin{proposition} $p: Y \twoheadrightarrow Y/G$ is regular and there is a short exact sequence
\[
1 \to \pi (Y, y_0) \xhookrightarrow{p_*} \pi (Y/G, x_0) \xtwoheadrightarrow{\varphi } G \to 1 
\] 
\end{proposition}

\begin{proposition} Let $p: \tilde{X} \twoheadrightarrow X$ covering space. Then if there exist a lifting $\tilde{\varphi }$, it induces a lifting on fundamental groups.
	% https://q.uiver.app/#q=WzAsOCxbMCwxLCJZIl0sWzEsMSwiWCJdLFsxLDAsIlxcdGlsZGV7WH0iXSxbMywxLCJcXHBpKFksIHlfMCkiXSxbNCwxLCJcXHBpKFgsIHhfMCkiXSxbNCwwLCJcXHBpKFxcdGlsZGV7WH0sIFxcdGlsZGV7eH1fMCkiXSxbMiwwLCJcXGltcGxpZXMiXSxbMiwxLCJcXGltcGxpZXMiXSxbMCwyLCJcXHRpbGRle1xcdmFycGhpfSJdLFswLDEsIlxcdmFycGhpIiwyXSxbMiwxLCJwIiwwLHsic3R5bGUiOnsiaGVhZCI6eyJuYW1lIjoiZXBpIn19fV0sWzMsNSwiXFx0aWxkZXtcXHZhcnBoaX1fKiJdLFszLDQsIlxcdmFycGhpIiwyXSxbNSw0LCJwXyoiLDAseyJzdHlsZSI6eyJoZWFkIjp7Im5hbWUiOiJlcGkifX19XV0=
\[\begin{tikzcd}
	& {\tilde{X}} & \implies && {\pi(\tilde{X}, \tilde{x}_0)} \\
	Y & X & \implies & {\pi(Y, y_0)} & {\pi(X, x_0)}
	\arrow["p", two heads, from=1-2, to=2-2]
	\arrow["{p_*}", two heads, from=1-5, to=2-5]
	\arrow["{\tilde{\varphi}}", from=2-1, to=1-2]
	\arrow["\varphi"', from=2-1, to=2-2]
	\arrow["{\tilde{\varphi}_*}", from=2-4, to=1-5]
	\arrow["\varphi"', from=2-4, to=2-5]
\end{tikzcd}\]
\end{proposition}

\begin{theorem} Let $p: \tilde{X} \twoheadrightarrow X$ covering space and $\tilde{x_0} \in p^{-1}(x_0)$. Then 
\[
	\exists \tilde{\varphi }: Y \to \tilde{X} \iff  \varphi _*(\pi (Y,y_0)) \subseteq p_*( \pi (\tilde{X}, \tilde{x}_0))
\]
\end{theorem}

\begin{proposition} If $Y$ is simply connected, then $\varphi : Y \to X $ has a lifting $\tilde{\varphi }: Y \to \tilde{X}$ (given $\pi  (Y, y_0) = 1$ )
\end{proposition}

\begin{definition}[Equivalence of covering spaces] $p_1: \tilde{X}_1 \to X$ and $p_2: \tilde{X}_2 \to X$ are equivalent if $\ \exists h: \tilde{X}_1 \xrightarrow{\cong} \tilde{X}_2$ homeomorphism such that
% https://q.uiver.app/#q=WzAsMyxbMCwwLCJcXHRpbGRle1h9XzEiXSxbMiwwLCJcXHRpbGRle1h9XzIiXSxbMSwxLCJYIl0sWzAsMiwicF8xIl0sWzEsMiwicF8yIiwyXSxbMCwxLCJcXGV4aXN0cyBoIl1d
\[\begin{tikzcd}
	{\tilde{X}_1} && {\tilde{X}_2} \\
	& X
	\arrow["{\exists h}", from=1-1, to=1-3]
	\arrow["{p_1}", from=1-1, to=2-2]
	\arrow["{p_2}"', from=1-3, to=2-2]
\end{tikzcd}\]
commutes
\end{definition}

\begin{theorem} $p_1: \tilde{X}_1 \to X$ and $p_2: \tilde{X}_2 \to X$ are equivalent if and only if $p_{1*}(\pi (\tilde{X}_1, \tilde{x}_1))$ and $p_{2*}(\pi (\tilde{X}_2, \tilde{x}_2))$ are conjugate.
\end{theorem}

\begin{theorem} Let $p: \tilde{X} \to X$ and $q: \tilde{X} \to \tilde{X}/A$ regular covering spaces. Then $X \cong \tilde{X} / A$.
\end{theorem}

\begin{proposition}
Regular covering spaces are all of the form $p: Y \to T/G$.
\end{proposition}

\begin{theorem}[Semilocally simply conntected]
$p: \tilde{X} \to X$ universal covering if and only if every point $x \in X$ admits a neighborhood $U$ such that $\pi (U, x) \to \pi (X, x)$ is trivial ($X$ semilocally simply connected).
\end{theorem}

\begin{theorem} Every subgroup of a free group is free. 
\end{theorem}













\end{document}